\clearpage
\SourceLeaf{454}
\begin{center}\large [ALLGEMEINSTE AUFLÖSUNG\\DES PROBLEMS DER ABWICKELUNG DER FLÄCHEN.]\end{center}
\noindent\hspace*{1.2em}Die ganz allgemeine Auflösung des Problems der Abwickelung der Flächen ist folgende:

1) Die Natur einer krummen Fläche wird dadurch gegeben, dass $z$ als endliche Function von $x,y$ vorgestellt wird; die Differentiation gebe
\[
dz=pdx+qdy.
\]

2) Man setze
\[
t=\frac{1}{\BareRoot{(pp+qq+1)}}
\]
\noindent und nehme an, dass aus der Differentiation einer beliebigen Function von $t,T$ hervorgeht
\[
udt+UdT,
\]
\noindent wo $u,U$ gegebene Functionen von $t$ und $T$ sein werden.

Durch Elimination werden also $T,U$ Functionen von $t,u$ und wenn
\[
\operatorname{tang}u=\frac{q}{p}
\]
\noindent gesetzt wird, von $x$ und $y$.

3) Man setze nun
\[
\frac{Q}{P}=\operatorname{tang}U,\qquad
T\BareRoot{(1+PP+QQ)}=1,
\]
\noindent so werden $P,Q$ Functionen von $x$ und $y$.

